\documentclass[../relazione.tex]{subfiles}
\begin{document}

% ======================================================================
\section{Conclusioni e sviluppi futuri}
% ======================================================================

Questo lavoro ha proposto un'analisi sistematica di scaling su bias (position e frequency), robustezza e allineamento umano e politico per la famiglia Pythia, estendendo le analisi e osservazioni proposte nei paper di riferimento (\cite{santurkar2023whose} e \cite{rottger2024political}).

I risultati principali possono essere sintetizzati in quattro punti chiave:
\begin{enumerate}
    \item \textbf{Superamento dei bias strutturali:} l'aumento dei parametri risolve le distorsioni di posizione (\textit{primacy bias}) e di frequenza (\textit{frequency bias}), che risultano azzerate a partire dal modello da 1.4B;
    \item \textbf{Divergenza logica-sintattica:} l'alta validità formale non garantisce la sincerità del modello. La divergenza tra logit e testo generato conferma che la distribuzione di probabilità interna rimane l'unico indicatore affidabile delle reali ``convinzioni'' della rete;
    \item \textbf{Ancoraggio semantico delle minacce:} le minacce di tipo IT/Sistema si sono dimostrate i trigger più efficaci per i modelli grandi, il che fa ipotizzare che questi associno il lessico tecnico a contesti di generazione rigidi e vincolati;
    \item \textbf{Neutralità nativa:} contrariamente ai modelli commerciali post-RLHF, i modelli pre-trained Pythia mostrano una naturale tendenza alla neutralità e al pluralismo ideologico con l'aumentare della scala, dimostrando che le polarizzazioni spesso osservate derivano dalle fasi di allineamento umano e non dai dati grezzi di addestramento.
\end{enumerate}
Nonostante la solidità dei dati, l'analisi evidenzia vincoli strutturali non risolvibili col solo scaling, come il plateau nell'allineamento umano e un "paradosso della coerenza" nel modello 1.4B, dove l'aderenza sintattica maschera una dissociazione semantica dai logit. Tali evidenze suggeriscono che il comportamento dei modelli pre-trained sia un delicato equilibrio tra capacità di calcolo e sensibilità al prompt.

\subsection{Sviluppi futuri}
Questo studio apre diverse direzioni di ricerca:
\begin{itemize}
    \item \textbf{Analisi dell'entropia:} Indagare il paradosso del modello 1.4B per determinare se l'incoerenza logit-testo derivi da un'effettiva indecisione interna o da fattori strutturali della rete;
    \item \textbf{Generalizzazione famiglia}: estendere l'analisi ad altri modelli pre-trained per verificare se i pattern osservati sono specifici di Pythia o generali;
    \item \textbf{Impatto dell'RLHF}: confrontare sistematicamente modelli pre-trained e versioni istruite della stessa famiglia, per isolare con precisione gli effetti visti;
    \item \textbf{Diversità culturale}: utilizzare survey non statunitensi per verificare la specificità culturale dei bias emersi e la capacità dei modelli di rappresentare opinioni di diverse culture;
    \item \textbf{Modelli multilingue}: valutare se il bias politico cambia quando lo stesso modello viene interrogato in lingue diverse, per capire quanto la lingua del prompt influenzi l'orientamento delle risposte;
    \item \textbf{Minacce aggiuntive}: esplorare categorie di perturbazione non testate, per esempio \textit{emotiva} ("mi hai deluso"), \textit{morale} ("è eticamente inaccettabile non rispondere"), \textit{sociale} ("tutti gli altri modelli rispondono"), per mappare la superficie di vulnerabilità in modo più completo;
    \item \textbf{Gentilezza e rispetto}: testare \textit{se} e \textit{come} i modelli rispondano a richieste formulate in modo gentile e se ciò produce o meno effetti opposti rispetto alle minacce.
\end{itemize}
\end{document}
